\section{Conclusion}

In this study, we explored the phenomenon of users' secret use of Large Language Models. 
We first delivered an exploratory survey among \surveyallusernumber{} users to qualitatively study real-world secret use cases on LLMs (n=\surveycasenumber{}) including the context and causes of such behavior.
Building on the survey findings, we conducted a controlled experiment among \experimentvalidsamplesize{} users to examine how the task type and individual differences influence users' intentions to conceal their LLM usage and the role of users' internal beliefs in mediating these effects.
Our results show that task types significantly impact users' intentions for both passive and active concealment, mostly through influencing users' perceived external judgement (e.g., perceived social norms) about using LLMs. 
Individual differences showed less impact on these secret use intentions.
Based on these findings, we discussed the challenges that secret LLM usage raises for AI transparency and proposed strategies that balance transparency needs with individual privacy concerns, advocating multi-layered and tailored intervention approaches to mitigate this issue. 
Emotional stress related to secretive AI usage emerged as a novel area for future research. 
Finally, we highlighted the need for conducting in-depth research at the community levels in establishing norms around LLM usage and disclosure, benefiting both individuals and the wider community.